\section{Penelitian Terkait}

Penelitian terkait ABAC untuk IoT umumnya menekankan kemampuan model berbasis atribut dalam menyediakan kontrol akses yang lebih \textit{fine-grained}, dinamis, dan \textit{context-aware} dibandingkan model tradisional. \textcite{diaz_authorization_2025} menjelaskan bahwa model otorisasi pada IoT menentukan aksi spesifik yang boleh dilakukan suatu entitas setelah proses autentikasi berlangsung. Sementara itu, \textcite{ullah_survey_2023} membahas integrasi \textit{blockchain} dengan ABAC sebagai pendekatan untuk meningkatkan keamanan dan kontrol akses pada lingkungan IoT. Namun, sebagian besar penelitian pada kelompok ini lebih banyak menekankan aspek ekspresivitas kebijakan, skalabilitas, desentralisasi, dan efisiensi mekanisme akses, sementara pembahasan mengenai \textit{freshness} atribut sebagai faktor yang dapat memengaruhi reliabilitas keputusan otorisasi masih relatif terbatas.

Kelompok penelitian otorisasi \textit{edge/fog} menunjukkan bahwa pemindahan proses pengambilan keputusan ke \textit{edge} atau \textit{fog} dapat mengurangi ketergantungan pada \textit{cloud} dan mendukung respons yang lebih cepat. \textcite{kalaria_adaptive_2024} mengusulkan kontrol akses adaptif berbasis \textit{fog} yang \textit{context-aware} dengan membawa proses pengambilan keputusan dan informasi kebijakan lebih dekat ke \textit{end nodes}. Namun, penurunan latensi tidak secara otomatis menjamin keputusan akses tetap akurat apabila atribut yang digunakan pada \textit{edge} sudah tidak mutakhir. \textcite{wang_edge-enabled_2025} secara eksplisit menyatakan bahwa \textit{edge-enabled IAM} yang mengasumsikan atribut selalu mutakhir dapat menghasilkan keputusan akses yang patut dipertanyakan pada lingkungan IoT yang dinamis dan terdistribusi.

Penelitian tentan sinkronisasi atribut dan revokasi mulai menyoroti pentingnya konsistensi atribut dan pencabutan hak akses. \textcite{wang_edge-enabled_2025} membahasa LAA di \textit{edge} dengan mempertimbangkan \textit{freshness} atribut dan \textit{trustworthiness} penyedia atribut. Selain itu \textcite{penuelas-angulo_revocation_2023} menegaskan bahwa revokasi merupakan permasalahan penting pada skema kontrol akses \textit{fine-grained} berbasis atribut. Namun, studi-studi tersebut belum secara eksplisit membahas optimasi konfigurasi TTL atribut, ukuran \textit{cache}, maupun pemilihan atribut yang layak di-\textit{cache} sebagai bagian dari mekanisme \textit{security-aware edge authorization}.

Penelitian \textit{edge caching} umumnya menunjukkan bahwa \textit{caching} dapat menurunkan latensi, meingkatkan responsivitas sistem, dan mengurangi biaya operasional pada lingkungan IoT. \textcite{weerasinghe_adaptive_2023} membahas \textit{adaptive context caching} untuk aplikasi IoT \textit{real-time}. Sementara, itu, \textcite{medvedev_refresh_2023} membahas \textit{trade-off} antara \textit{freshness}, latensi, QoS, dan biaya operasional dalam strategi \textit{caching} dan \textit{prefetching} pada IoT \textit{middleware}. Namun, fokus utama penelitian tersebut masih berada pada \textit{context/data caching}, bukan \textit{attribute caching} untuk pengambilan keputusan keamanan berbasis ABAC.

Penelitian metaheuristik pada IoT dan \textit{fog computing} juga menunjukkan efektivitas algoritma seperti GWO untuk penjadwalan tugas dan alokasi sumber daya. \textcite{alsadie_modified_2025} menggunakan GWO yang dimodifikasi untuk penjadwalan tugas pada \textit{fog/cloud computing}, sedangkan \textcite{alsamarai_improved_2024} menggunakan PC-GWO untuk meningkatkan performa penjadwalan, utilisasi sumber daya, dan efisiensi energi pada lingkungan \textit{fog-cloud}. Akan tetapi, penggunaan GWO pada penelitian tersebut masih difokuskan pada optimasi performa komputasi, bukan optimasi parameter keamanan ABAC seperti TTL atribut, ukuran \textit{cache}, dan \textit{cacheable attribute}.

Berdasarkan tinjauan literatur tersebut, masih terbatas penelitian yang secara terpadu menggabungkan \textit{edge-based ABAC, attribute caching, freshness-aware authorization}, optimasi berbasis metaheuristik, dan evaluasi berorientasi keamanan terhadap kualitas keputusan akses. Oleh karena itu, penelitian ini mengevaluasi konfigurasi ABAC berbasis \textit{edge} menggunakan metrik performa, \textit{overhead, freshness}, dan kualitas keputusan secara bersamaan. Dengan demmikian, sistem tidak hanya dinilai dari sisi latensi atau \textit{cache hit rate}, tetapi juga dari kemampuannya menjaga reliabilitas dan keamanan keputusan akses ketika atribut berubah secara dinamis.

\begin{landscape}
\begin{table}[H]
\centering
\scriptsize
\caption{Perbandingan penelitian terkait ABAC, caching, dan optimasi pada IoT}
\setlength{\tabcolsep}{2pt}
\renewcommand{\arraystretch}{1.25}

\begin{tabularx}{\linewidth}{|
L{2.4cm}|
L{2.1cm}|
L{2.2cm}|
L{1.2cm}|
L{1.2cm}|
L{2.4cm}|
L{2.0cm}|
L{2.8cm}|
Y|}
\hline
\textbf{Penelitian} &
\textbf{Kelompok} &
\textbf{Domain} &
\textbf{Edge/Fog} &
\textbf{ABAC} &
\textbf{Cache/Freshness} &
\textbf{Optimasi} &
\textbf{Metrik/Evaluasi} &
\textbf{Keterbatasan} \\ \hline

\textcite{diaz_authorization_2025}
& ABAC/ authorization IoT
& IoT authorization survey
& Tidak utama
& Ya
& Terbatas
& Tidak
& Analisis komparatif model dan arsitektur otorisasi
& Membahas authorization IoT dan kebutuhan \textit{fine-grained/context-aware}, tetapi freshness atribut belum menjadi fokus utama. \\ \hline

\textcite{ullah_survey_2023}
& ABAC untuk IoT
& Blockchain-based ABAC for IoT
& Tidak utama
& Ya
& Terbatas
& Tidak
& Security/privacy, latency, overhead, scalability
& Fokus pada ABAC dan blockchain; belum membahas \textit{stale attribute} dan optimasi TTL/cache. \\ \hline

\textcite{kalaria_adaptive_2024}
& Edge/fog authorization
& Fog-based context-aware access control
& Ya
& Ya
& Tidak utama
& Tidak
& Processing time overhead, security improvement
& Fog menurunkan latency dan cloud dependency, tetapi belum fokus pada \textit{stale attribute}/freshness atribut. \\ \hline

\textcite{wang_edge-enabled_2025}
& Edge-IAM; synchronization
& Edge-enabled IAM IoT
& Ya
& Ya
& Attribute freshness dan synchronization
& Tidak
& Freshness, source trustworthiness, decision reliability, overhead
& Membahas freshness dan local authorization, tetapi belum mengoptimalkan TTL, cache size, dan cacheable attributes. \\ \hline

\textcite{penuelas-angulo_revocation_2023}
& Revocation berbasis atribut
& ABE revocation for fog-enabled IoT
& Ya
& Tidak langsung
& Revocation, bukan cache freshness
& Tidak
& User revocation, attribute revocation, cost comparison
& Membahas revocation, tetapi belum mengaitkannya dengan optimasi TTL/cache atribut untuk edge ABAC. \\ \hline

\textcite{weerasinghe_adaptive_2023}
& Context caching
& IoT context caching/CMP
& Tidak utama
& Tidak
& Context freshness
& Adaptive/RL-based caching
& Cost efficiency, QoS/QoC, freshness, latency, hit rate
& Fokus pada context/data caching, bukan cache atribut keamanan ABAC. \\ \hline

\textcite{medvedev_refresh_2023}
& IoT caching
& IoT middleware/CMP caching
& Tidak utama
& Tidak
& Freshness--latency--cost
& Cost-based caching/prefetching
& QoS, freshness, latency, operational cost
& Membahas trade-off caching, tetapi bukan otorisasi ABAC. \\ \hline

\textcite{alsadie_modified_2025}
& Metaheuristik IoT
& Fog-cloud task scheduling
& Ya
& Tidak
& Tidak
& Modified GWO
& Makespan, energy consumption
& GWO efektif untuk scheduling, tetapi belum diarahkan ke optimasi parameter ABAC. \\ \hline

\textcite{alsamarai_improved_2024}
& Metaheuristik IoT
& Fog-cloud task scheduling
& Ya
& Tidak
& Tidak
& PC-GWO
& Makespan, energy, resource utilization, cost
& Optimasi performa/resource allocation, bukan security-aware ABAC. \\ \hline

Penelitian ini
& ABAC edge security-aware optimization
& Edge IoT authorization
& Ya
& Ya
& TTL per atribut, stale hit
& Random Search, GA, GWO
& FAR, FDR, stale hit, revocation FAR, latency, PIP calls
& Simulasi terkontrol; belum diuji pada deployment fisik. \\ \hline

\end{tabularx}
\label{tab:related-work-abac-iot}
\end{table}
\end{landscape}